%% ============================================================
%%  Chapter 4 — Results
%% ============================================================
\chapter{Results}
\label{ch:results}

This chapter reports quantitative outcomes from eight sets of experiments:
(1) SEDM cross-intent evaluation,
(2) per-intent action distribution analysis,
(3) composite risk and escalation risk characterisation,
(4) DQN training dynamics,
(5) parameter-selection and robustness analysis,
(6) the effect of the traffic-realism extensions on classifier accuracy and
classifier-driven false positives,
(7) window vs.\ EMA escalation-tracking comparison, and
(8) cross-session reputation and bounded threshold-adaptation behaviour.
Section~\ref{sec:results_sedm}--Section~\ref{sec:results_params} report
results for the system as validated prior to this version — these numbers
are unchanged and are retained as the baseline against which
Section~\ref{sec:results_realism}--Section~\ref{sec:results_adaptive}
measure the effect of the new extensions.
All results were obtained using the protocols and hyperparameters described
in Chapter~\ref{ch:methodology}.

\paragraph{Methodology note (evaluation alignment).} Every result in this
chapter scores each defender action against the ground-truth label of the
\emph{same} observation the action was chosen from, decoded directly from
the state vector before the environment is advanced. An earlier version of
the evaluation pipeline paired each action with the \emph{following} step's
label instead — a one-step lag that silently corrupts false-positive and
detection-rate figures while leaving reward and action-distribution figures
largely unaffected (documented as Bug 8/9 in \texttt{docs/BUGS\_AND\_FIXES.md}).
All numbers below reflect the corrected, re-run evaluation.

%% ── 4.1 SEDM Cross-Intent Evaluation ────────────────────────
\section{SEDM Cross-Intent Evaluation}
\label{sec:results_sedm}

\subsection{Overall Performance Summary}

Table~\ref{tab:eval_summary} summarises SEDM performance across all four
attacker intents, evaluated over 30 episodes of 200 steps each, with the
SEDM deciding from the RandomForest classifier's \emph{predicted} attack
type (classifier-driven mode — see Section~\ref{sec:results_extended} and
Section~\ref{sec:results_params} for the oracle-vs-classifier distinction
and the robustness analysis behind the near-zero false positive rate).
The SEDM achieves detection rates exceeding 99.9\% across all four profiles,
demonstrating robust and consistent threat coverage without any
intent-specific training or tuning.

\begin{table}[ht]
  \centering
  \caption{SEDM evaluation summary across all attacker intents
           (30 episodes $\times$ 200 steps each, classifier-driven decisions).
           Mean $\pm$ standard deviation are reported.}
  \label{tab:eval_summary}
  \begin{tabular}{lcccccc}
    \toprule
    Intent & Mean Reward & Std Reward & Det.\ Rate & FP Rate & Avg Threat & Avg Risk \\
    \midrule
    STEALTHY      & 1010.22 & 51.52 & \textbf{99.93\%} & 0.00\% & 0.806 & 0.632 \\
    AGGRESSIVE    & 1090.84 & 19.80 & \textbf{100.0\%} & 0.00\% & 0.854 & 0.709 \\
    TARGETED      & 1126.10 & 21.54 & \textbf{99.97\%} & 0.00\% & 0.853 & 0.673 \\
    OPPORTUNISTIC &  895.05 & 33.30 & \textbf{99.97\%} & 0.00\% & 0.790 & 0.662 \\
    \bottomrule
  \end{tabular}
\end{table}

\subsection{Detection Rate Analysis}

Detection rates are uniformly high across all four intent profiles,
ranging from 99.93\% (STEALTHY) to 100.0\% (AGGRESSIVE).
The sub-0.1 percentage-point spread across profiles indicates that the
SEDM's intent-aware escalation risk computation adapts consistently across
structurally different attack campaigns, and that the residual misses are
a classifier artefact rather than a policy weakness: the classifier's
held-out accuracy is 99.85\% (10-class), with its only measurable confusion
between EXPLOITS/GENERIC and their structurally similar neighbours
(Section~\ref{sec:results_params}).

Under the STEALTHY intent, the campaign proceeds slowly through early kill
chain stages with low escalation risk.
The SEDM assigns Low-band actions (ALLOW, LOG) at early stages, correctly
identifying benign-like reconnaissance.
When STEALTHY campaigns eventually progress to later stages (Installation,
C2), the escalation risk increases, and the SEDM appropriately escalates
to BLOCK/ALERT.
The small residual miss rate (0.07\%) traces to the rare cases where the
classifier misreads a genuine low-severity attack sample as NORMAL,
triggering the R1 override.

The AGGRESSIVE intent achieves the highest detection rate (100.0\% across
30 evaluation episodes) because its high escalation rate and fast kill-chain
progression place almost every step in the High escalation-risk band, where
the matrix lookup returns BLOCK or ALERT regardless of the exact attack type
— making the policy's response robust to residual classifier noise on this
particular intent.

\subsection{False Positive Rate Analysis}

Measured false positive rates are 0.00\% for all four intents under the
classifier-driven evaluation. This number should not be read as a claim
that the SEDM is immune to false positives in general — it is the
consequence of two specific, checked properties of the current experimental
setup, not an emergent property of the kill-chain policy itself:

\begin{enumerate}
  \item The R1 override unconditionally maps a predicted-NORMAL sample to
        ALLOW, so \emph{any} false positive must originate upstream, in the
        classifier mistaking a true NORMAL sample for an attack type.
  \item The RandomForest classifier achieves 100\% precision and recall on
        the NORMAL class specifically (99.85\% accuracy overall), because
        the synthetic feature simulator (\texttt{attacker/attack\_types.py})
        draws each attack type from a disjoint parametric distribution by
        construction.
\end{enumerate}

Because real network telemetry does not offer this clean separability, a
feature-noise robustness sweep was run to check how the false-positive rate
degrades as the classifier's input signal becomes noisier — see
Section~\ref{sec:results_params}. The sweep recovers a non-zero,
monotonically increasing false-positive rate (up to 10.00\% at 50\%
injected feature noise), confirming the zero-FP headline number is a
property of the current synthetic simulator's separability, not evidence
that the SEDM would be immune to false positives against real, noisy
traffic. Section~\ref{sec:results_realism} shows that this same
NORMAL-recall ceiling is also the property that shifts, in a small but
measurable way, once the traffic-realism extensions of
Section~\ref{sec:realism} are introduced.

\subsection{Reward Distribution}

Mean episode rewards range from 895.05 (OPPORTUNISTIC) to 1126.10 (TARGETED).
The reward is determined by the alignment between the SEDM's actions and
the reward matrix (Table~\ref{tab:reward_matrix_meth}).

TARGETED produces the highest reward because its concentrated,
high-severity, late-stage attacks are consistently met with BLOCK/ALERT,
which yield the highest rewards for high-threat traffic.
Standard deviations are low (19--21 for TARGETED and AGGRESSIVE),
reflecting the deterministic SEDM policy applied to relatively homogeneous
intent-specific campaigns.

OPPORTUNISTIC yields the lowest mean reward (895.05) with higher standard
deviation (33.30).
The scattered attack-type distribution in OPPORTUNISTIC campaigns results
in episodes where medium-severity attacks (FUZZERS, GENERIC at Delivery
stage) receive TROLL responses rather than BLOCK, incurring suboptimal
rewards.
The higher episode-to-episode variance reflects the greater diversity of
attack sequences generated by the OPPORTUNISTIC Markov chain.

STEALTHY yields moderate mean reward (1010.22) with the highest standard
deviation (51.52), consistent with its inherently variable behaviour:
some STEALTHY episodes remain in early kill chain stages for extended
periods (producing many LOG/ALLOW rewards), while others escalate to
late stages (producing BLOCK/ALERT rewards).
This bimodal episode structure drives the elevated standard deviation.

\begin{figure}[ht]
  \centering
  \includegraphics[width=0.82\textwidth]{reward_boxplot}
  \caption{Episode reward distributions per attacker intent (30 episodes each).
           Boxes show the interquartile range; whiskers extend to 1.5$\times$IQR;
           circles indicate outliers.}
  \label{fig:reward_boxplot}
\end{figure}

\subsection{Revised Primary Protocol: Mixed Traffic and Confidence Intervals}
\label{sec:results_primary_revised}

The false-positive rate reported above (0.00\% for all four intents) is
scored against whichever NORMAL-labelled samples the attacker's own Markov
chain happened to produce during the \texttt{benign\_ratio}$=0$ protocol —
there is no explicit benign-traffic injection. Counting these samples
directly exposes how small that denominator actually is:
\textbf{33--114 benign samples per intent} (Table~\ref{tab:primary_ci_comparison}),
out of 6\,000 total steps. A point estimate computed from so few samples of
the relevant class carries very little statistical weight on its own,
independent of whether the policy or classifier is well-calibrated.

To make this explicit rather than leaving it implicit,
Table~\ref{tab:primary_ci_comparison} re-runs the identical classifier-driven
decision logic used to produce Table~\ref{tab:eval_summary} (same seed, same
\texttt{MatrixPolicy}, same predicted-vs-ground-truth scoring) under two
protocols side by side: the original protocol as already reported, and a
revised primary protocol with \texttt{benign\_ratio}$=0.3$, using the
50-episode $\times$ 500-step budget already justified for mixed-traffic
conditions in Table~\ref{tab:sem_episodes} (Section~\ref{sec:results_params}).
Both protocols report detection rate and false-positive rate as Wilson score
95\% confidence intervals computed from the raw TP/FP/TN/FN counts, rather
than as bare percentages, specifically so a rate estimated from a handful of
samples is visibly less certain than one estimated from thousands.

\begin{table}[htbp]
  \centering
  \small
  \caption{Detection rate and false-positive rate under the original
           (\texttt{benign\_ratio}$=0.0$, 30$\times$200) and revised
           (\texttt{benign\_ratio}$=0.3$, 50$\times$500) primary protocols,
           with Wilson 95\% confidence intervals and raw sample counts
           (classifier-driven decisions, identical decision logic to
           Table~\ref{tab:eval_summary}).}
  \label{tab:primary_ci_comparison}
  \begin{tabular}{llrlrl}
    \toprule
    Intent & Protocol & $n$(attack) & Det.\ Rate [95\% CI] & $n$(benign) & FP Rate [95\% CI] \\
    \midrule
    STEALTHY      & Original & 5\,886  & 99.98\% [99.90, 99.997] & 114   & 0.00\% [0.00, 3.26]  \\
    STEALTHY      & Revised  & 17\,251 & 99.95\% [99.91, 99.98]  & 7\,749 & 0.68\% [0.52, 0.89]  \\
    AGGRESSIVE    & Original & 5\,960  & 99.98\% [99.91, 99.997] & 40    & 0.00\% [0.00, 8.76]  \\
    AGGRESSIVE    & Revised  & 17\,445 & 99.94\% [99.89, 99.96]  & 7\,555 & 0.60\% [0.45, 0.80]  \\
    TARGETED      & Original & 5\,967  & 99.92\% [99.80, 99.96]  & 33    & 0.00\% [0.00, 10.43] \\
    TARGETED      & Revised  & 17\,446 & 99.94\% [99.89, 99.96]  & 7\,554 & 0.79\% [0.62, 1.02]  \\
    OPPORTUNISTIC & Original & 5\,956  & 100.0\% [99.94, 100.0]  & 44    & 0.00\% [0.00, 8.03]  \\
    OPPORTUNISTIC & Revised  & 17\,429 & 99.97\% [99.93, 99.99]  & 7\,571 & 0.65\% [0.49, 0.85]  \\
    \bottomrule
  \end{tabular}
\end{table}

\textbf{The original protocol's 0.00\% figure was correct but barely
informative.} Its confidence interval upper bound ranges from 3.26
percentage points (STEALTHY, 114 benign samples) up to 10.43 percentage
points (TARGETED, only 33 benign samples). In other words, a measured
0.00\% false-positive rate under the original protocol was statistically
consistent with a true underlying false-positive rate as high as roughly
1-in-10 for TARGETED — the point estimate did not rule out a materially
non-zero rate, it simply reflected that too few benign samples existed to
observe one.

\textbf{The revised protocol multiplies the benign sample count by
70--230$\times$} (e.g.\ 33 $\to$ 7\,554 for TARGETED, 114 $\to$ 7\,749 for
STEALTHY), narrowing each false-positive-rate confidence interval to
roughly 0.2--0.4 percentage points wide, and reveals a small, real, and now
precisely-bounded false-positive rate of \textbf{0.60--0.79\%} across all
four intents. This is not a contradiction of the 0.00\% figure reported for
the original protocol — the two numbers measure the same policy and the
same classifier under different traffic compositions, and the entire
difference between them is explained by the change in benign sample count,
not by any change in SEDM or classifier behaviour. Detection rate, by
contrast, is essentially unchanged between protocols (99.94--99.97\%
vs.\ 99.92--100.0\%, with heavily overlapping confidence intervals),
confirming that mixed-traffic injection specifically sharpens the
false-positive estimate without altering detection performance.

\textbf{Independent agreement with Table~\ref{tab:realism_mixed_metrics}.}
The 0.60--0.79\% false-positive rate measured here, using
\texttt{models/classifier.joblib} and a purpose-built scoring script, is
directly consistent with the specificity figures already reported in
Table~\ref{tab:realism_mixed_metrics} (0.992--0.995, i.e.\ an implied
0.5--0.8\% false-positive rate) computed independently via the
\texttt{evaluation.sedm\_eval} infrastructure under the same
50$\times$500, \texttt{benign\_ratio}$=0.3$ protocol. Two independently
implemented measurement paths agreeing on the same magnitude is stronger
evidence for the underlying number than either measurement alone.

\textbf{Adopted as the primary protocol going forward.} Because it is both
more statistically rigorous (adequate sample sizes for both classes under
evaluation) and more operationally realistic (real deployments are not
continuous, uninterrupted attack traffic), the mixed-traffic,
confidence-interval-reported protocol in Table~\ref{tab:primary_ci_comparison}
is the standard this thesis recommends for future primary evaluations,
without retracting Table~\ref{tab:eval_summary}: Table~\ref{tab:eval_summary}'s
numbers remain valid as a measurement of the SEDM under continuous-attack
conditions, they are simply not, on their own, adequate evidence about
false-positive behaviour under realistic mixed traffic — a distinction this
section makes explicit rather than leaving implicit.

\begin{figure}[ht]
  \centering
  \includegraphics[width=0.88\textwidth]{metric_comparison_revised_primary}
  \caption{Detection rate and false-positive rate under the revised primary
           protocol (50 episodes $\times$ 500 steps, 30\% benign-traffic
           injection), directly comparable to Table~\ref{tab:eval_summary}'s
           continuous-attack figures.}
  \label{fig:metric_comparison_revised_primary}
\end{figure}

%% ── 4.2 Action Distribution Analysis ────────────────────────
\section{Per-Intent Action Distribution}
\label{sec:results_action}

Table~\ref{tab:action_dist} reports the percentage of steps assigned to each
honeypot action across the four attacker intents.

\begin{table}[ht]
  \centering
  \caption{Action distribution across attacker intents (all 30 evaluation
           episodes pooled, classifier-driven decisions).
           Values are percentages of total steps.}
  \label{tab:action_dist}
  \begin{tabular}{lrrrrr}
    \toprule
    Intent & ALLOW & LOG & TROLL & BLOCK & ALERT \\
    \midrule
    STEALTHY      & 2.0\% & 0.3\% & 0.9\% & 1.7\% & 95.1\% \\
    AGGRESSIVE    & 0.7\% & 0.0\% & 0.1\% & 1.7\% & 97.6\% \\
    TARGETED      & 0.6\% & 0.0\% & 0.1\% & 1.5\% & 97.8\% \\
    OPPORTUNISTIC & 0.8\% & 0.0\% & 0.2\% & 2.8\% & 96.2\% \\
    \bottomrule
  \end{tabular}
\end{table}

\subsection{Dominance of ALERT}

ALERT alone accounts for 95.1--97.8\% of actions across all intent profiles,
with BLOCK a distant second (1.5--2.8\%). This is more heavily ALERT-skewed
than the pre-fix figures reported in an earlier draft of this evaluation
(which showed a wider BLOCK/ALERT split); the corrected label alignment
(Section~\ref{ch:results}, methodology note) removes a spurious source of
BLOCK attributions that were artefacts of the one-step lag rather than
genuine policy behaviour. It remains consistent with the evaluation design:
all four intents involve continuous attack traffic with no extended benign
periods (this is the 30-episode$\times$200-step, \texttt{benign\_ratio}$=0$
protocol — see Section~\ref{sec:results_extended} for the mixed-traffic
variant), keeping the threat level elevated and the kill-chain stage at
advanced positions for most of the episode.

The near-total absence of ALLOW (0.6--2.0\%) confirms that the SEDM
correctly avoids passthrough for ongoing attack campaigns; the small
residual reflects genuine NORMAL-type traffic naturally sampled by the
attacker's own Markov chain, not misclassification
(Section~\ref{sec:results_params}).

\subsection{AGGRESSIVE and TARGETED: Near-Total ALERT}

AGGRESSIVE (97.6\% ALERT) and TARGETED (97.8\% ALERT) are almost entirely
ALERT responses.
Both intents produce high escalation risk throughout their campaigns:
AGGRESSIVE because of strong forward kill chain transitions;
TARGETED because of the concentrated exploitation-to-installation path.
High escalation risk places most stages in the High band, mapping to
ALERT at Installation, C2, and Actions on Objectives (Table~\ref{tab:sedm_matrix}).
The R2 override (DOS/WORMS $\to$ upgrade) further increases ALERT frequency
for AGGRESSIVE campaigns, which generate significant DOS and WORMS traffic.

\subsection{STEALTHY and OPPORTUNISTIC: Slightly More Graded}

STEALTHY (95.1\% ALERT, 2.0\% ALLOW, 1.7\% BLOCK) and OPPORTUNISTIC (96.2\%
ALERT, 2.8\% BLOCK) retain the largest non-ALERT fractions of the four
intents, consistent with their design: STEALTHY spends more of the episode
at early kill-chain stages with Low/Medium escalation risk before
committing to an attack path, and OPPORTUNISTIC's scattered attack-type
distribution produces more Medium-band DELIVERY/EXPLOITATION steps that map
to TROLL/BLOCK rather than ALERT. Both effects are the same mechanism
identified in Section~\ref{sec:results_risk} (escalation risk by kill-chain
stage), now visible with the corrected label alignment rather than
partially masked by it.

\begin{figure}[ht]
  \centering
  \includegraphics[width=0.90\textwidth]{effective_policy_per_intent}
  \caption{Stacked bar chart of action distribution across the four attacker
           intents. Each bar represents 30 pooled evaluation episodes
           (6\,000 steps total per intent).}
  \label{fig:action_dist}
\end{figure}

%% ── 4.3 Composite Risk Characterisation ──────────────────────
\section{Escalation and Composite Risk Analysis}
\label{sec:results_risk}

\subsection{Composite Risk Score Distribution}

The composite risk score $\rho_c$ (Section~\ref{sec:sedm_alg}, Step 5)
provides an aggregate threat index at each step.
Mean composite risk scores range from 0.632 (STEALTHY) to 0.709 (AGGRESSIVE),
consistent with the attack severity ordering of the four intents.

AGGRESSIVE campaigns achieve the highest mean composite risk (0.709) because
they combine high attack severity (DOS, WORMS, EXPLOITS) with high escalation
probability and high escalation rate.
STEALTHY campaigns have the lowest mean composite risk (0.632) despite having
the highest detection rate variability; their low escalation probability and
lower attack severities (RECONNAISSANCE, ANALYSIS) reduce the composite score.

\begin{figure}[ht]
  \centering
  \includegraphics[width=0.82\textwidth]{composite_risk_distribution}
  \caption{Distribution of composite risk scores across all evaluation
           steps, stratified by attacker intent.
           Higher risk scores correspond to more dangerous attack configurations.}
  \label{fig:composite_risk}
\end{figure}

\subsection{Escalation Risk by Kill Chain Stage}

Escalation risk $\rho(k, \pi)$ as a function of kill chain stage for each
intent profile (Figure~\ref{fig:esc_risk_per_intent}) is computed
analytically from the intent-specific Markov transition matrices.
Several patterns are noteworthy:
%
\begin{itemize}
  \item For AGGRESSIVE, escalation risk is consistently high ($>0.65$)
        from WEAPONIZATION onwards, explaining the dominance of High-band
        matrix entries and the near-total ALERT response.
  \item For STEALTHY, escalation risk is low ($<0.35$) at early stages
        (RECONNAISSANCE, WEAPONIZATION) and rises only at Installation
        and beyond.
        This explains why the SEDM assigns ALLOW/LOG at early STEALTHY
        stages and only escalates to BLOCK/ALERT when the campaign
        reaches Installation.
  \item For TARGETED, escalation risk is high from EXPLOITATION onwards,
        consistent with the focused exploitation chain.
        Early stages (RECONNAISSANCE, WEAPONIZATION) have moderate risk,
        reflecting the targeted campaign's tendency to skip early stages
        or pass through them quickly.
  \item For OPPORTUNISTIC, escalation risk is moderate ($0.35$--$0.65$)
        across most stages, placing many steps in the Medium band.
        This drives the significant BLOCK fraction (24.6\%) as Medium-band
        EXPLOITATION steps map to BLOCK.
\end{itemize}

\begin{figure}[ht]
  \centering
  \includegraphics[width=0.90\textwidth]{escalation_risk_per_intent}
  \caption{Escalation risk $\rho(k, \pi)$ as a function of kill chain stage,
           for all four attacker intent profiles.
           Dashed horizontal lines show the Low (0.35) and High (0.65)
           band thresholds.}
  \label{fig:esc_risk_per_intent}
\end{figure}

\subsection{SEDM Decision Matrix Visualisation}

A colour-coded representation of the $7 \times 3$ SEDM
(Figure~\ref{fig:sedm_heatmap}), with action severity encoded on a
five-point colour scale (ALLOW = lightest, ALERT = darkest), shows a clear
escalation gradient from top-left (Reconnaissance/Low = ALLOW) to
bottom-right (Actions on Objectives/High = ALERT), confirming the
proportionality principle underlying the matrix design.

\begin{figure}[ht]
  \centering
  \includegraphics[width=0.65\textwidth]{sedm_decision_matrix}
  \caption{Stage-Escalation Decision Matrix (SEDM) visualised as a
           colour-coded heatmap.
           Rows represent kill chain stages (Reconnaissance at top,
           Actions on Objectives at bottom);
           columns represent escalation risk bands
           (Low on left, High on right).}
  \label{fig:sedm_heatmap}
\end{figure}

%% ── 4.4 Kill Chain Stage Distribution ───────────────────────
\section{Kill Chain Stage Distribution}
\label{sec:results_killchain}

\begin{figure}[ht]
  \centering
  \includegraphics[width=0.88\textwidth]{kill_chain_distribution}
  \caption{Distribution of kill chain stages visited across evaluation
           episodes, stratified by attacker intent.
           Bars show the fraction of steps in each stage.}
  \label{fig:killchain_dist}
\end{figure}

The kill chain stage distribution across evaluation episodes
(Figure~\ref{fig:killchain_dist}) for each intent is consistent with the
intent descriptions in Section~\ref{sec:intents}:
%
\begin{itemize}
  \item \textbf{STEALTHY}: overrepresented at RECONNAISSANCE and
        WEAPONIZATION; relatively rare at ACTIONS\_ON\_OBJ.
  \item \textbf{AGGRESSIVE}: concentrated at EXPLOITATION, INSTALLATION,
        and COMMAND\_AND\_CTRL; rapid progression means less time in
        early stages.
  \item \textbf{TARGETED}: concentrated at EXPLOITATION and INSTALLATION,
        consistent with the focused exploit chain.
  \item \textbf{OPPORTUNISTIC}: roughly uniform across mid-chain stages
        (DELIVERY through COMMAND\_AND\_CTRL), reflecting the scattered
        transition structure.
\end{itemize}

%% ── 4.5 DQN Training Dynamics ────────────────────────────────
\section{DQN Training Dynamics}
\label{sec:results_dqn}

\subsection{Overview}

The DQN baseline was trained for 300 episodes of 500 steps each on the
OPPORTUNISTIC intent.
Training metrics were logged to \texttt{logs/metrics.csv} and visualised
as six-panel training curves.

\begin{figure}[ht]
  \centering
  \includegraphics[width=0.92\textwidth]{training_curves}
  \caption{DQN training curves over 300 episodes.
           \emph{Top row, left to right}: episode total reward (raw and
           rolling 10-episode mean); episode detection rate; episode
           false-positive rate.
           \emph{Bottom row, left to right}: average episode threat level;
           per-update Huber loss (smoothed); per-episode average loss.}
  \label{fig:training_curves}
\end{figure}

\subsection{Episode Reward}

Episode reward climbs rapidly from 1\,006 at episode~0 to approximately
2\,051 at episode~1, then continues to rise to a plateau of 2\,200--2\,360
by episode~30 (Figure~\ref{fig:training_curves}, top-left).
The dramatic improvement between episodes 0 and 1 reflects the structure of
the OPPORTUNISTIC attacker: even a partially trained policy quickly learns
that non-ALLOW actions yield positive rewards on the dense attack traffic,
causing a substantial reward increase as soon as the replay buffer contains
enough transitions for a meaningful Q-value update.

The rolling 10-episode mean stabilises at approximately 2\,300 from
episode~50 onwards, with relatively small episode-to-episode variance
($\sigma \approx 35$), indicating that the policy has converged to a
stable strategy.
The absence of reward collapse (a common failure mode in DQN training)
suggests that the target network and gradient clipping together provide
adequate stabilisation.

\subsection{Detection Rate}

Detection rate starts at 87.6\% at episode~0, rises sharply to 97.4\% at
episode~1, and exceeds 98.5\% from episode~3 onwards
(Figure~\ref{fig:training_curves}, top-centre).
The high starting detection rate (87.6\%) is consistent with the
OPPORTUNISTIC attacker's continuous attack traffic: even a random policy
that selects non-ALLOW actions for the majority of steps will detect most
attacks because the benign fraction of traffic is small.

Detection rates fluctuate within a narrow band ($98.5$--$99.5\%$) throughout
training, with occasional dips to $\approx$\,97\% corresponding to episodes
where the epsilon-greedy policy selects suboptimal ALLOW responses more
frequently than usual.

\subsection{False-Positive Rate}

The false-positive rate exhibits a distinctive pattern: it remains at
essentially 1.0 throughout most of training, with occasional drops to 0.0
at isolated episodes (e.g., episodes 62 and 142).
This bimodal behaviour reflects the fact that false positive events are
rare in the OPPORTUNISTIC scenario (the attacker generates predominantly
attack traffic), so the per-episode false-positive rate is either 1.0
(the episode happened to contain a benign step and the DQN responded
with a non-ALLOW action) or 0.0 (no benign steps occurred, making the
false positive rate undefined or effectively zero).

This finding reveals an important distinction between the DQN training
environment and the SEDM evaluation: the DQN training set is dominated
by attack traffic, providing limited signal for learning to discriminate
benign traffic.
The SEDM's explicit ALLOW entry in the matrix for Reconnaissance/Low
provides a structural guarantee that benign-like traffic at early kill
chain stages is handled correctly, without requiring the learner to observe
rare benign examples.

\subsection{DQN Training Loss}

Per-update Huber loss increases from $\approx$\,0.97 at training start
(episode~0) to $\approx$\,1.4--1.6 by episodes 1--10, then gradually
rises to $\approx$\,2.0--2.6 and stabilises with moderate variance through
episodes 50--300 (Figure~\ref{fig:training_curves}, bottom panels).

The rising loss during early training is expected: as the policy network
learns to produce larger Q-value estimates for highly rewarding actions, the
Huber targets also grow, causing the raw loss value to increase even as the
\emph{relative} TD error decreases.
The absence of large upward spikes in the later training phase indicates
that gradient clipping and the target network together prevent instability.
The persistent moderate variance in per-update loss ($\approx$\,$\pm$0.5$\sigma$
around the rolling mean) is characteristic of off-policy DQN training with a
diverse replay buffer.

\begin{table}[ht]
  \centering
  \caption{DQN training performance at selected training milestones
           (OPPORTUNISTIC intent, 500 steps per episode).}
  \label{tab:dqn_milestones}
  \begin{tabular}{rrrrrr}
    \toprule
    Episode & Total Reward & Det.\ Rate & FP Rate & Avg Threat & Avg Loss \\
    \midrule
    0   & 1\,006.0 & 87.6\% & 100.0\% & 0.808 & 0.967 \\
    1   & 2\,051.6 & 97.4\% & 100.0\% & 0.808 & 1.415 \\
    10  & 2\,288.0 & 99.0\% & 100.0\% & 0.808 & 1.661 \\
    50  & 2\,285.0 & 99.2\% & 100.0\% & 0.808 & 2.207 \\
    100 & 2\,303.5 & 99.0\% & 100.0\% & 0.808 & 2.404 \\
    150 & 2\,297.1 & 99.0\% & 100.0\% & 0.808 & 2.442 \\
    200 & 2\,217.4 & 98.2\% & 100.0\% & 0.808 & 1.942 \\
    250 & 2\,356.4 & 99.6\% & 100.0\% & 0.808 & 2.133 \\
    299 & 2\,272.4 & 98.8\% & 100.0\% & 0.808 & 2.453 \\
    \bottomrule
  \end{tabular}
\end{table}

\subsection{Comparison with SEDM}

The DQN and SEDM are compared on the metrics available from both systems.
On the OPPORTUNISTIC intent (the DQN's training distribution), the DQN
achieves detection rates of 98.5--99.6\% with training reward
2\,200--2\,360 per 500-step episode.
The SEDM achieves 99.41\% detection on OPPORTUNISTIC in 200-step evaluation
episodes.

However, the DQN's false-positive rate is effectively undefined in training
(benign steps are too rare for a meaningful estimate), while the SEDM
achieves 15.00\% on OPPORTUNISTIC in evaluation episodes that include more
diverse traffic.
The DQN requires 300 training episodes (150\,000 environment steps) to
reach its performance level; the SEDM requires no training.

The DQN's opaque Q-values provide no direct insight into why particular
actions are chosen, while the SEDM's five-step algorithm is fully
traceable to observable inputs.
This interpretability advantage is examined further in Chapter~\ref{ch:discussion}
(Discussion), and is the specific evidence Section~5.3 relies on when
arguing against a learning-based mechanism for the dynamic-response
extensions in Section~\ref{sec:dynamic_response}.

%% ── 4.6 Metric Comparison Summary ───────────────────────────
\section{Cross-Intent Metric Comparison}
\label{sec:results_compare}

Figure~\ref{fig:metric_comparison} presents a comparative view of the four
key metrics across attacker intents.

\begin{figure}[ht]
  \centering
  \includegraphics[width=0.88\textwidth]{metric_comparison}
  \caption{Side-by-side comparison of mean reward, detection rate,
           false positive rate, and average threat level for each attacker
           intent.
           Error bars show $\pm$1 standard deviation over 30 episodes.}
  \label{fig:metric_comparison}
\end{figure}

Figure~\ref{fig:radar_comparison} shows the same metrics as a radar chart,
enabling simultaneous comparison across all four dimensions.
The radar chart reveals that TARGETED achieves the best overall balance:
highest detection rate, lowest false positive rate, and highest mean reward.
STEALTHY, while maintaining $>$\,99\% detection, exhibits a larger false
positive radius, reflecting the challenge of distinguishing reconnaissance
traffic from benign traffic at low escalation risk.

\begin{figure}[ht]
  \centering
  \includegraphics[width=0.70\textwidth]{radar_comparison}
  \caption{Radar chart comparing SEDM performance across attacker intents
           on four normalised metrics:
           detection rate, false positive rate (inverted: lower is better),
           mean reward (normalised), and average threat level.}
  \label{fig:radar_comparison}
\end{figure}

%% ── 4.7 Extended SEDM Metrics ───────────────────────────────────────
\section{Extended SEDM Metrics: Precision, Proportionality, and Containment}
\label{sec:results_extended}

Section~\ref{sec:results_sedm}--Section~\ref{sec:results_compare} characterise
SEDM performance under scenarios composed entirely of attack traffic
(30~episodes $\times$ 200~steps, no benign traffic).
This section reports a supplementary evaluation conducted under \emph{mixed}
traffic conditions to provide a more realistic operational picture:
50~episodes of 500~steps each per intent, with a 30\% benign-traffic injection
ratio (\texttt{benign\_ratio}$= 0.30$).
The additional metrics reported here are precision, recall, F\textsubscript{1},
F\textsubscript{2}, specificity, response proportionality, late-stage miss rate,
and mean steps-to-containment.

\subsection{Classification Metrics Under Mixed Traffic}

Table~\ref{tab:sedm_extended_metrics} reports the full set of binary
classification metrics for the SEDM, treating any non-ALLOW response as a
positive detection, for both evaluation modes: \textbf{oracle} (the SEDM
decides from ground-truth kill-chain state — an upper bound on achievable
performance) and \textbf{classifier-driven} (the SEDM decides from the
RandomForest's predicted attack type — the realistic operating condition,
since ground truth is never directly observable in deployment).

\begin{table}[ht]
  \centering
  \caption{Extended classification metrics for the SEDM policy across all
           four attacker intents (50~episodes $\times$ 500~steps each,
           benign-traffic ratio~30\%), oracle vs.\ classifier-driven decisions.
           Mean values over all evaluation episodes are reported.}
  \label{tab:sedm_extended_metrics}
  \begin{tabular}{llcccccc}
    \toprule
    Intent & Variant & Precision & Recall & F\textsubscript{1} & F\textsubscript{2} & Specificity & Prop.\ Score \\
    \midrule
    STEALTHY      & Oracle     & 1.000 & 1.000 & 1.000 & 1.000 & 1.000 & 0.693 \\
    STEALTHY      & Classifier & 1.000 & 1.000 & 1.000 & 1.000 & 0.999 & 0.693 \\
    AGGRESSIVE    & Oracle     & 1.000 & 1.000 & 1.000 & 1.000 & 1.000 & 0.698 \\
    AGGRESSIVE    & Classifier & 1.000 & 1.000 & 1.000 & 1.000 & 0.999 & 0.698 \\
    TARGETED      & Oracle     & 1.000 & 1.000 & 1.000 & 1.000 & 1.000 & 0.699 \\
    TARGETED      & Classifier & 1.000 & 0.999 & 0.999 & 0.999 & 0.999 & 0.699 \\
    OPPORTUNISTIC & Oracle     & 1.000 & 1.000 & 1.000 & 1.000 & 1.000 & 0.698 \\
    OPPORTUNISTIC & Classifier & 0.999 & 1.000 & 1.000 & 1.000 & 0.999 & 0.698 \\
    \bottomrule
  \end{tabular}
\end{table}

\paragraph{Precision, recall, and specificity.}
Under the oracle variant, precision, recall, F\textsubscript{1}, F\textsubscript{2}, and
specificity are exactly 1.000 for all four intents. This is provable rather
than merely observed: the R1 override maps every ground-truth-NORMAL state
to ALLOW and every other state falls in a matrix band that is never ALLOW
for the stages reachable under these four intents, so the oracle SEDM cannot
produce a false positive or false negative by construction (this is the
same tautology documented as a follow-up to Bug 8 in
\texttt{docs/BUGS\_AND\_FIXES.md}).

The classifier-driven variant is the meaningful test of practical
performance: it breaks the tautology by deciding from the RandomForest's
\emph{predicted} attack type while still scoring against ground truth.
Results remain at 0.999--1.000 across all metrics because the classifier
itself is 99.85\% accurate on held-out data (Section~\ref{sec:results_params}),
with essentially perfect separation of the NORMAL class specifically. The
0.1 percentage-point gap between oracle and classifier-driven TARGETED
recall (1.000 vs.\ 0.999) is the one place this ceiling is visible in the
extended evaluation, and traces directly to the classifier's residual
EXPLOITS/GENERIC confusion. Section~\ref{sec:results_realism} reports the
same table's response to the traffic-realism extensions of
Section~\ref{sec:realism}.

\paragraph{Reading this table together with Section~\ref{sec:results_params}.}
Because both variants report near-1.000 metrics, this table alone cannot
distinguish "the SEDM is a good policy" from "the evaluation is too easy."
The feature-noise robustness sweep in Section~\ref{sec:results_params} is
the control that resolves this: it shows the classifier-driven numbers
degrade smoothly and non-trivially once the input signal is made
realistically noisy, which is the evidence that the near-perfect
classifier-driven results above are a genuine (if optimistic) measurement,
not a methodological artefact.

\begin{figure}[ht]
  \centering
  \includegraphics[width=0.88\textwidth]{sedm_classification_metrics}
  \caption{Precision, Recall, F\textsubscript{1}, and F\textsubscript{2}
           for each attacker intent under the mixed-traffic evaluation
           (50~episodes, 30\% benign ratio).}
  \label{fig:sedm_class_metrics}
\end{figure}

\paragraph{Response proportionality score.}
The proportionality score measures the fraction of steps at which the SEDM
assigns an action at least as severe as the minimum expected for that kill
chain stage (Table~\ref{tab:min_severity}).
Scores of 0.690--0.695 indicate that the SEDM meets the minimum severity
threshold at approximately 69\% of all steps.
The remaining 31\% correspond predominantly to early-stage
(RECONNAISSANCE, WEAPONIZATION) attack steps at which the SEDM issues
ALLOW by design (matrix entry [RECON, Low]~=~ALLOW), while the
minimum expected severity is defined as LOG.
This is an intentional design property of the SEDM: passive observation
at the earliest kill chain stages maximises threat intelligence without
prematurely disrupting the attacker's campaign.

\begin{table}[ht]
  \centering
  \caption{Minimum expected action severity per kill chain stage, used for
           the proportionality score computation.
           Action severities: ALLOW\,=\,0, LOG\,=\,1, TROLL\,=\,2,
           BLOCK\,=\,3, ALERT\,=\,4.}
  \label{tab:min_severity}
  \begin{tabular}{lcc}
    \toprule
    Kill Chain Stage & Min.\ Expected Severity & Min.\ Action \\
    \midrule
    RECONNAISSANCE   & 0 & ALLOW \\
    WEAPONIZATION    & 1 & LOG   \\
    DELIVERY         & 1 & LOG   \\
    EXPLOITATION     & 2 & TROLL \\
    INSTALLATION     & 3 & BLOCK \\
    COMMAND \& CTRL  & 3 & BLOCK \\
    ACTIONS ON OBJ.  & 4 & ALERT \\
    \bottomrule
  \end{tabular}
\end{table}

\subsection{Stage-wise Action Distribution}

Table~\ref{tab:stage_action_dist} reports the per-stage action distribution
pooled across all four intents and all evaluation episodes,
restricted to genuine attack steps (benign steps are excluded; they always
receive ALLOW via the R1 override regardless of stage).

\begin{table}[ht]
  \centering
  \caption{Per-stage action distribution for attack steps, pooled
           across all four attacker intents and all 50 evaluation episodes
           per intent (classifier-driven decisions; benign steps excluded,
           they always receive ALLOW via the R1 override).
           Values are row-normalised percentages; each row sums to 100\%.}
  \label{tab:stage_action_dist}
  \begin{tabular}{lrrrrr}
    \toprule
    Kill Chain Stage & ALLOW & LOG & TROLL & BLOCK & ALERT \\
    \midrule
    Reconnaissance   & 0.0\% & 40.0\% & 60.0\% &  0.0\% &  0.0\% \\
    Weaponization    & 0.0\% & 18.4\% & 43.2\% & 38.3\% &  0.0\% \\
    Delivery         & 0.0\% &  0.0\% & 69.2\% & 30.8\% &  0.0\% \\
    Exploitation     & 0.0\% &  0.0\% &  0.0\% & 73.0\% & 27.0\% \\
    Installation     & 0.1\% &  0.0\% &  0.0\% & 28.0\% & 71.9\% \\
    Command \& Ctrl  & 0.0\% &  0.0\% &  0.0\% &  0.0\% & 100.0\% \\
    Actions on Obj.  & 0.0\% &  0.0\% &  0.0\% &  0.0\% & 100.0\% \\
    \bottomrule
  \end{tabular}
\end{table}

The corrected table shows a clean, monotonic escalation ladder with no
residual ALLOW column at late stages: LOG and TROLL are confined to
RECONNAISSANCE/WEAPONIZATION/DELIVERY, BLOCK peaks at EXPLOITATION/
INSTALLATION (73.0\%/28.0\%), and ALERT reaches 100\% at COMMAND\,\&\,CONTROL
and ACTIONS ON OBJECTIVES. An earlier draft of this table (produced before
the label-alignment fix, Section~\ref{ch:results}) showed a persistent
$\sim$30\% ALLOW column at every late stage and was narrated as attacker
"camouflage traffic"; that pattern has not reproduced under the corrected
evaluation (the single residual 0.1\% ALLOW cell at INSTALLATION is
consistent with ordinary classifier noise, not a systematic camouflage
signal) and the camouflage interpretation is withdrawn — it was an artefact
of scoring actions against the wrong step's ground truth, not a property of
attacker behaviour or the SEDM's R1 rule.

\begin{figure}[ht]
  \centering
  \includegraphics[width=0.88\textwidth]{sedm_stage_action_heatmap}
  \caption{Row-normalised stage$\times$action distribution for attack steps,
           pooled across all intents and evaluation episodes.
           Warmer colours indicate higher frequency;
           the escalation gradient from top-left (RECON/ALLOW) to
           bottom-right (AOO/ALERT) is clearly visible.}
  \label{fig:stage_action_heatmap}
\end{figure}

A Spearman rank correlation of $\rho = 0.492$ ($p < 10^{-200}$) between kill
chain stage index and action severity, computed over all attack steps under
the classifier-driven variant, confirms a statistically significant,
moderate-to-strong positive monotonic relationship — more than four times
the $\rho = 0.115$ measured under the pre-fix (misaligned) evaluation. The
increase is expected: the earlier measurement mixed each action with the
\emph{next} step's stage/severity pairing, diluting the true stage-severity
relationship with essentially random noise. $\rho$ is well below 1.0
because RECONNAISSANCE/WEAPONIZATION intentionally use LOG/TROLL rather
than the numerically lowest action (ALLOW is reserved for confirmed-benign
traffic, Table~\ref{tab:sedm_matrix}), compressing the low end of the
severity range relative to a strictly linear stage-to-severity mapping.
Within the stages where the SEDM is expected to escalate (EXPLOITATION
through ACTIONS ON OBJ.), BLOCK and ALERT together account for 100\% of
steps at every one of those four stages (Table~\ref{tab:stage_action_dist})
— a clean, complete escalation trend, not merely a statistically
significant one.

\begin{figure}[ht]
  \centering
  \includegraphics[width=0.88\textwidth]{sedm_proportionality}
  \caption{Mean action severity per kill chain stage (orange circles) against
           the minimum expected severity (grey squares).
           The SEDM meets or exceeds the minimum at all stages from
           DELIVERY onwards; the gap at RECONNAISSANCE and WEAPONIZATION
           reflects the deliberate ALLOW policy for early-stage traffic.}
  \label{fig:sedm_proportionality}
\end{figure}

\subsection{Threat Containment Metrics}

Table~\ref{tab:sedm_containment} quantifies two containment-specific metrics:
the \emph{late-stage miss rate} and the \emph{mean steps to containment},
under the classifier-driven variant.

\begin{table}[ht]
  \centering
  \caption{Threat containment metrics per attacker intent
           (50~episodes $\times$ 500~steps, benign ratio 30\%,
           classifier-driven decisions).
           \emph{Late-stage miss rate} is the fraction of
           Installation$+$ attack steps receiving a weak response
           (ALLOW or LOG).
           \emph{Steps-to-contain} is the mean steps from the first attack
           step to the first BLOCK or ALERT response.}
  \label{tab:sedm_containment}
  \begin{tabular}{lcccc}
    \toprule
    Intent & Late-Stage Miss & Steps-to-Contain & Avg Threat & Avg Risk \\
    \midrule
    STEALTHY      & 0.000\% & 2.3 & 0.678 & 0.636 \\
    AGGRESSIVE    & 0.000\% & 0.2 & 0.708 & 0.665 \\
    TARGETED      & 0.001\% & 0.6 & 0.708 & 0.665 \\
    OPPORTUNISTIC & 0.000\% & 1.1 & 0.665 & 0.622 \\
    \bottomrule
  \end{tabular}
\end{table}

\paragraph{Late-stage miss rate.}
The late-stage miss rate is effectively 0\% across all four intents under
the corrected evaluation, not the $\approx$30\% reported in an earlier draft
of this chapter. The 30\% figure was an artefact of the one-step label lag
(Section~\ref{ch:results}, methodology note): actions chosen in response to a
late-stage attack step were being scored against the \emph{following} step's
ground truth, and roughly 30\% of the time that following step happened to
be a different (often lower-severity) sample, registering as a "miss" that
never actually occurred. With the alignment corrected, the "camouflaging
attacker" interpretation offered in an earlier draft is withdrawn — the
SEDM does not exhibit a systematic blind spot to NORMAL-signature traffic
at late kill-chain stages in this evaluation. Session-level camouflage
(an attacker deliberately alternating attack and benign-looking traffic
across \emph{multiple} sessions to build trust before an attack) remains a
theoretically valid concern and is retained as a limitation in
Chapter~\ref{ch:discussion}, but the empirical support for it claimed here
previously does not survive the bug fix and should not be cited as evidence
for that concern. The cross-session reputation mechanism introduced in
Section~\ref{sec:reputation} is, notably, a direct structural response to
this same concern — see Section~\ref{sec:results_reputation} and
Section~5.4.

\paragraph{Steps-to-containment.}
Steps-to-containment measures the response latency from the first attack step
to the first BLOCK or ALERT.
Values range from 0.2~steps (AGGRESSIVE) to 2.3~steps (STEALTHY) — faster
across the board than the pre-fix figures (1.24--3.22~steps), because the
corrected alignment removes the artificial one-step delay baked into every
prior measurement.

AGGRESSIVE campaigns achieve near-immediate containment (0.2 steps) because
their high escalation rate and early High-band classification trigger
BLOCK/ALERT from the first attack step. STEALTHY campaigns require more
steps (2.3) because early-stage STEALTHY traffic falls in the Low or Medium
escalation band, producing LOG/TROLL responses until the escalation risk
crosses the High threshold — this is intentional SEDM behaviour
(Table~\ref{tab:sedm_matrix}), not a detection delay.

\begin{figure}[ht]
  \centering
  \includegraphics[width=0.88\textwidth]{sedm_containment}
  \caption{Left: Late-stage miss rate by attacker intent (now $\approx$0\%
           after the label-alignment fix).
           Right: Mean steps-to-containment by attacker intent.
           AGGRESSIVE achieves near-immediate containment (0.2~steps);
           STEALTHY requires the most steps (2.3) before the escalation
           risk crosses the High-band threshold.}
  \label{fig:sedm_containment}
\end{figure}

\paragraph{Summary.}
The extended evaluation confirms that the SEDM delivers consistent,
intent-independent classification performance (F\textsubscript{1}
$\ge 0.999$ under classifier-driven decisions) and near-instantaneous
containment (0.2--2.3~steps) across all four attacker profiles, with no
measurable late-stage miss rate. The genuinely open limitation is not a
policy blind spot but an evaluation-realism one: both the near-1.000
classification metrics here and the 0\% headline false-positive rate in
Section~\ref{sec:results_sedm} follow from the synthetic feature simulator's
clean class separability, quantified precisely by the classifier's 99.85\%
held-out accuracy. Section~\ref{sec:results_params} shows this ceiling is
not a free assumption — accuracy and false-positive rate degrade
predictably as feature-measurement noise increases — but validating
end-to-end SEDM performance against real (non-synthetic) network captures
remains outside the scope of this evaluation and is identified as future
work in Chapter~\ref{ch:discussion}. Section~\ref{sec:results_realism}
reports the first concrete step toward closing this gap from within the
simulation itself: making the synthetic traffic generator less trivially
separable.

%% ── 4.8 Parameter Selection and Robustness Analysis ─────────────────
\section{Parameter Selection and Robustness Analysis}
\label{sec:results_params}

The preceding sections report point estimates for a fixed set of
protocol parameters (episode/step counts, escalation-risk thresholds,
classifier hyperparameters, benign-traffic ratio). This section documents
the empirical and design rationale behind each choice, so the reported
metrics can be read as the output of a justified protocol rather than
arbitrarily-chosen settings. Full detail and reproduction commands are in
\texttt{docs/parameter\_selection.md} and \texttt{docs01/parameter\_selection.md}.

\subsection{Episode and Step Counts}

The primary evaluation protocol (30 episodes $\times$ 200 steps) and the
extended protocol (50 episodes $\times$ 500 steps) were chosen using the
running standard error of the mean episode reward as a stopping criterion,
computed over 80 rollouts per intent:

\begin{table}[ht]
  \centering
  \caption{Standard error of the mean episode reward (as a percentage of
           the mean) as a function of episode count, for the two intents
           with the highest inter-episode variance.}
  \label{tab:sem_episodes}
  \begin{tabular}{lcccccc}
    \toprule
    Intent & $n$=5 & $n$=10 & $n$=20 & $n$=30 & $n$=50 & $n$=80 \\
    \midrule
    STEALTHY      & 1.08\% & 2.02\% & 1.18\% & 0.94\% & 0.72\% & 0.57\% \\
    OPPORTUNISTIC & 1.46\% & 1.47\% & 0.90\% & 0.68\% & 0.61\% & 0.48\% \\
    \bottomrule
  \end{tabular}
\end{table}

Both intents fall below 1\% relative SEM by $n=30$, with diminishing returns
beyond that point ($n=50 \to n=80$ improves SEM by only $\sim$0.15
percentage points). 30 episodes was therefore adopted as the primary
evaluation budget; 50 episodes for the extended (mixed-traffic) evaluation
tightens the interval further given the additional variance introduced by
benign-traffic injection.

\subsection{Escalation Risk Bands and Override Thresholds}

\texttt{ESC\_LOW\_THRESHOLD}$=0.35$ and \texttt{ESC\_HIGH\_THRESHOLD}$=0.65$
(Table~\ref{tab:sedm_matrix}) were chosen as an equal three-way split of the
$[0,1]$ escalation-risk probability space, rather than fit to any one
intent's transition matrix — fitting the thresholds to the four intents
used in evaluation would risk overstating the cross-intent generalisation
claimed in Section~\ref{sec:results_sedm}. Because the SEDM is a fixed
lookup table, its sensitivity to these thresholds is fully auditable: a
$\pm 0.05$ perturbation changes the assigned band for at most one adjacent
stage per intent, and Figure~\ref{fig:esc_risk_per_intent} shows the
escalation-risk values for AGGRESSIVE and TARGETED sit well clear of both
boundaries, while STEALTHY's early stages sit closest to the 0.35 boundary
— consistent with STEALTHY being the lowest-margin intent on every metric
reported in this chapter. \texttt{RATE\_THRESHOLD}$=0.80$ (the R3 override
trigger) is set near the top of the escalation-rate range so it fires only
under sustained, dense attack traffic, avoiding spurious escalation on
ordinary bursty traffic. As Section~\ref{sec:results_escalation_mode} and
Section~\ref{sec:results_adaptive} show, this design intent is only partly
realised under continuous-attack conditions with the original window-based
escalation signal, where \texttt{RATE\_THRESHOLD} is in practice exceeded
for the majority of an episode's duration.

\subsection{Classifier Hyperparameters}

The RandomForest (\texttt{n\_estimators=150}, \texttt{max\_depth=20},
\texttt{class\_weight=\textquotesingle balanced\textquotesingle}) reaches 99.85\% held-out
accuracy at these settings. Because the underlying feature distributions
(\texttt{attacker/attack\_types.py}) are well-separated parametric
distributions by construction, classifier capacity is not the limiting
factor here, so these were kept as conventional, non-overfit defaults
rather than tuned further — a grid search against this data would optimise
noise rather than signal. These hyperparameters are unchanged in
Section~\ref{sec:results_realism}'s retrained classifier; only the training
data generator changed.

\subsection{Feature-Noise Robustness Sweep}

The near-zero false-positive rate reported throughout this chapter follows
directly from the classifier's near-perfect separation of NORMAL from
attack traffic on the synthetic feature simulator. To check whether this
result is an artefact of unrealistically clean synthetic data rather than a
property of the SEDM policy, the classifier was re-evaluated with
multiplicative Gaussian noise injected onto every continuous feature before
prediction:

\begin{table}[ht]
  \centering
  \caption{Classifier accuracy and NORMAL$\to$Attack false-positive rate
           under increasing injected feature-measurement noise
           ($\sigma$, multiplicative Gaussian, applied per-feature).}
  \label{tab:noise_robustness}
  \begin{tabular}{lcc}
    \toprule
    Feature noise ($\sigma$) & Accuracy & NORMAL$\to$Attack FPR \\
    \midrule
    0\%  & 99.85\% & 0.00\%  \\
    5\%  & 99.60\% & 0.00\%  \\
    10\% & 99.35\% & 0.00\%  \\
    20\% & 97.75\% & 1.00\%  \\
    35\% & 94.60\% & 7.00\%  \\
    50\% & 87.65\% & 10.00\% \\
    \bottomrule
  \end{tabular}
\end{table}

\begin{figure}[ht]
  \centering
  \includegraphics[width=0.75\textwidth]{classifier_noise_robustness}
  \caption{Classifier accuracy (left axis) and NORMAL$\to$Attack false
           positive rate (right axis) as a function of injected feature
           measurement noise. False positives emerge smoothly once feature
           fidelity degrades beyond $\sim$10\% noise.}
  \label{fig:noise_robustness}
\end{figure}

The classifier is robust to small measurement noise ($\le$10\%, plausible
for well-instrumented flow export), and false positives emerge smoothly and
monotonically as feature fidelity degrades further, reaching 10.00\% FPR at
50\% noise. This is the direct evidence that the 0\% headline false-positive
rate is a property of the current synthetic simulator's clean separability,
not a general claim about deployment conditions: \textbf{the SEDM/classifier
pipeline does produce false positives under realistically degraded input,
exactly as expected of any statistical classifier}, and the magnitude is
now quantified rather than asserted. Section~\ref{sec:results_realism}
shows that traffic-realism alone — with no artificial noise injection at
all — moves the classifier in the same direction as this sweep, for the
same underlying reason (reduced class separability), by construction
rather than by injected noise.

\subsection{Data Format Generality}

The 15-field feature schema (\texttt{FEATURE\_NAMES}) is named after
UNSW-NB15 for familiarity, but neither the classifier nor the SEDM has a
code-level dependency on UNSW-NB15 or NSL-KDD specifically — both operate
on a plain \texttt{dict[str, float]} matching that schema. The same fields
map onto other commonly-available flow-export formats (NetFlow v9/IPFIX,
Zeek \texttt{conn.log}, CICFlowMeter output) via a straightforward field
mapping (see \texttt{docs/parameter\_selection.md}), so the architecture is
not bound to a single benchmark dataset's schema. This directly addresses
the practicality concern raised of prior, dataset-locked evaluation work:
the contribution under test is the kill-chain-aware SEDM policy and its
evaluation methodology, not a fixed dataset binding — though substituting a
real log parser for the synthetic simulator, and re-running this chapter's
evaluation against it, remains future work (Chapter~\ref{ch:discussion}).

%% ── 4.9 Effect of Traffic-Realism Extensions ─────────────────────────
\section{Effect of Traffic-Realism Extensions}
\label{sec:results_realism}

This section quantifies the effect of the session-coherent intensity/persona
extension (Section~\ref{sec:realism}) on the classifier and on the primary
SEDM evaluation protocol of Section~\ref{sec:results_sedm}. The classifier
was retrained from scratch (\texttt{fit\_from\_simulation}) on the extended
generator, using the identical hyperparameters and training-set size as
Section~\ref{sec:results_sedm} (6\,000 samples, 600 per class, seed 42), and
re-evaluated under the identical protocol used to produce the 99.85\%
figure reported in Section~\ref{sec:bg_rf}/Section~\ref{sec:results_params}
(2\,000 held-out samples, 200 per class, seed 999), so the two accuracy
figures are directly comparable.

\begin{table}[ht]
  \centering
  \small
  \caption{Classifier accuracy before and after the traffic-realism
           extensions of Section~\ref{sec:realism}, identical evaluation
           protocol (200 held-out samples per class, seed 999).}
  \label{tab:realism_accuracy}
  \begin{tabular}{lccc}
    \toprule
    Generator & Accuracy & NORMAL Precision & NORMAL Recall \\
    \midrule
    Original (\S\ref{sec:bg_rf}, \S\ref{sec:results_sedm}) & 99.85\% & 1.0000 & 1.0000 \\
    Session-coherent + personas (this section) & 99.40\% & 1.0000 & 0.9900 \\
    \bottomrule
  \end{tabular}
\end{table}

The 0.45 percentage-point accuracy drop is modest but not zero, and traces
to a specific, interpretable cause rather than a generalised degradation:
per-class F1 scores show the residual confusion concentrated in the same
EXPLOITS/GENERIC pair already identified in Section~\ref{sec:results_sedm}
as the classifier's weakest boundary (F1 = 0.977 and 0.975 respectively,
down from figures at or near 1.000), and — more directly relevant to the
false-positive story of Section~\ref{sec:results_sedm} — NORMAL-class
recall is no longer a perfect 1.0000 but 0.9900. Two effects plausibly
compound here: session-level intensity scaling widens the within-class
variance of the volume-shaped features for every attack type including
NORMAL, and the introduction of three distinct benign personas
(Section~\ref{sec:realism}) means the NORMAL class itself is no longer a
single, tight distribution but a mixture of three, some of which
(\texttt{crawler} in particular) share feature ranges with low-severity
attack types such as RECONNAISSANCE. Both are the intended effect of the
extension — a classifier trained and evaluated on more realistic, less
artificially separable traffic — not a regression.

\paragraph{Primary protocol re-evaluation.}
Re-running the exact Section~\ref{sec:results_sedm} protocol (30 episodes
$\times$ 200 steps per intent, classifier-driven, seed 42) with the
retrained classifier and the session-coherent generator reproduces
detection rates and mean rewards within the range attributable to ordinary
seed-level variation (detection rates 99.92--100.0\%, compared to
99.93--100.0\% in Table~\ref{tab:eval_summary}), and the measured
false-positive rate remains 0.00\% for all four intents in this specific
protocol. This last point requires an honest qualification rather than
being read as "no change": the primary protocol injects no explicit benign
traffic (\texttt{benign\_ratio}$=0$), so the only NORMAL-labelled samples it
contains are those the attacker's own Markov chain happens to produce,
which is too small a sample for a 1-percentage-point recall change to
register as a nonzero measured FPR. The "Revised Primary Protocol"
subsection at the end of Section~\ref{sec:results_sedm}
(Section~\ref{sec:results_primary_revised}) addresses this directly by
making mixed-traffic evaluation, with sample counts and confidence
intervals, the standard primary protocol; the table below is the
mixed-traffic instrument already established in
Section~\ref{sec:results_extended}, applied here as the more sensitive test
of the same question.

\paragraph{Mixed-traffic re-evaluation (the sensitive instrument).}
Table~\ref{tab:realism_mixed_metrics} re-runs the exact
Section~\ref{sec:results_extended} extended-metrics protocol (50 episodes
$\times$ 500 steps, 30\% benign-traffic ratio, classifier-driven) with the
retrained, realism-extended classifier, directly comparable to
Table~\ref{tab:sedm_extended_metrics}.

\begin{table}[ht]
  \centering
  \caption{Extended classification metrics under the traffic-realism
           extensions of Section~\ref{sec:realism}, identical protocol to
           Table~\ref{tab:sedm_extended_metrics} (50 episodes $\times$ 500
           steps, benign ratio 30\%, classifier-driven decisions).}
  \label{tab:realism_mixed_metrics}
  \begin{tabular}{lcccccc}
    \toprule
    Intent & Precision & Recall & F\textsubscript{1} & F\textsubscript{2} & Specificity & Prop.\ Score \\
    \midrule
    STEALTHY      & 0.997 & 0.999 & 0.998 & 0.999 & 0.993 & 0.695 \\
    AGGRESSIVE    & 0.998 & 0.999 & 0.999 & 0.999 & 0.995 & 0.699 \\
    TARGETED      & 0.997 & 0.999 & 0.998 & 0.999 & 0.993 & 0.700 \\
    OPPORTUNISTIC & 0.997 & 0.999 & 0.998 & 0.999 & 0.992 & 0.700 \\
    \bottomrule
  \end{tabular}
\end{table}

Here — where the protocol has enough true-negative samples to detect it —
the effect is unambiguous and consistent across all four intents:
specificity drops from 0.999 (Table~\ref{tab:sedm_extended_metrics}) to
0.992--0.995, and precision drops from 0.999--1.000 to 0.997--0.998. This is
the same underlying degradation Table~\ref{tab:realism_accuracy} measures
directly on the classifier in isolation, now visible end-to-end in the full
SEDM decision pipeline, on a protocol with adequate statistical power to
detect it (a Spearman correlation of $\rho = 0.490$ between kill-chain stage
and action severity, $p < 10^{-200}$, essentially unchanged from
Table~\ref{tab:sedm_extended_metrics}'s $\rho = 0.492$, confirms the core
stage-escalation relationship the SEDM is built around is undisturbed —
only the classifier-driven false-positive margin moved).

\paragraph{Interpretation.}
The result should be read as the traffic-realism extensions successfully
doing what they were designed to do: making the classification problem
measurably, if modestly, harder in a way that is directly traceable to
specific, intended changes (session coherence, persona diversity) rather
than to an implementation defect, and — crucially — a way that only becomes
\emph{visible} once the evaluation protocol itself has enough benign
samples to reveal it (see Section~\ref{sec:results_primary_revised}). The
zero-FPR headline figure in Section~\ref{sec:results_sedm} and the
near-perfect classification metrics in Table~\ref{tab:sedm_extended_metrics}
were always explicitly qualified in this thesis as consequences of the
synthetic simulator's separability rather than general robustness claims
(Section~\ref{sec:results_sedm}, Section~\ref{sec:results_params}); this
section and that subsection together provide the first concrete,
protocol-adequate evidence, generated without any artificial noise
injection, of that separability beginning to erode under a more realistic
generator.

%% ── 4.10 Escalation Tracking: Window vs. EMA ─────────────────────────
\section{Escalation Tracking: Window vs.\ EMA}
\label{sec:results_escalation_mode}

This section compares the window-based and severity-weighted-EMA escalation
signals (Section~\ref{sec:escalation_tracking}) directly, using the SEDM in
oracle mode (deciding from ground-truth state) so that the comparison
isolates the effect of the escalation signal itself from classifier noise.
For each of the four attacker intents, 30 episodes of 200 steps were run
under both \texttt{escalation\_mode} settings, using the same seed and
continuous-attack (\texttt{benign\_ratio}$=0$) protocol as
Section~\ref{sec:results_sedm}.

\begin{table}[htbp]
  \centering
  \small
  \caption{R3 trigger rate and escalation-signal statistics under window
           vs.\ EMA escalation tracking (30 episodes $\times$ 200 steps per
           intent, oracle mode, continuous attack traffic).}
  \label{tab:escalation_mode}
  \begin{tabular}{llcccccc}
    \toprule
    Intent & Mode & Det.\ Rate & FP Rate & R3 Trigger Rate & Mean $r_t$ & 95th pct.\ $r_t$ & Max $r_t$ \\
    \midrule
    STEALTHY      & window & 1.000 & 0.000 & \textbf{90.0\%} & 0.973 & 1.000 & 1.000 \\
    STEALTHY      & EMA    & 1.000 & 0.000 & \textbf{2.8\%}  & 0.676 & 0.799 & 0.835 \\
    AGGRESSIVE    & window & 1.000 & 0.000 & \textbf{44.7\%} & 0.991 & 1.000 & 1.000 \\
    AGGRESSIVE    & EMA    & 1.000 & 0.000 & \textbf{8.0\%}  & 0.755 & 0.849 & 0.875 \\
    TARGETED      & window & 1.000 & 0.000 & \textbf{67.7\%} & 0.994 & 1.000 & 1.000 \\
    TARGETED      & EMA    & 1.000 & 0.000 & \textbf{22.0\%} & 0.761 & 0.841 & 0.879 \\
    OPPORTUNISTIC & window & 1.000 & 0.000 & \textbf{68.6\%} & 0.990 & 1.000 & 1.000 \\
    OPPORTUNISTIC & EMA    & 1.000 & 0.000 & \textbf{0.7\%}  & 0.633 & 0.797 & 0.860 \\
    \bottomrule
  \end{tabular}
\end{table}

\begin{figure}[ht]
  \centering
  \includegraphics[width=0.82\textwidth]{escalation_mode_r3_rate}
  \caption{R3 override trigger rate by intent under window vs.\
           severity-weighted-EMA escalation tracking, showing a large and
           consistent reduction under EMA tracking.}
  \label{fig:escalation_mode_r3_rate}
\end{figure}

\textbf{Detection rate and false-positive rate are unaffected.} Both
signals produce identical oracle-mode detection (100\%) and false-positive
(0\%) rates for all four intents; this is expected under oracle decisions,
since R1's unconditional NORMAL$\to$ALLOW mapping does not depend on the
escalation signal at all, and the matrix-lookup component of the decision
(Table~\ref{tab:sedm_matrix}) depends on $\rho(k, \pi)$, not $r_t$. The
escalation signal's only influence on the final action, at oracle-mode
precision, is through R3.

\textbf{A large, consistent gap in R3 trigger rate.} Under the window
signal, R3 fires on 44.7--90.0\% of all decisions across the four intents —
under the continuous-attack protocol used here, the window-based escalation
rate saturates near its ceiling of 1.0 within the first $\sim$20 steps of an
episode and remains there for essentially the rest of the episode, so
\texttt{RATE\_THRESHOLD} (0.80) is exceeded for the majority of most
episodes' duration. Under the EMA signal, R3 fires far more selectively —
0.7--22.0\% — because the signal's own ceiling under sustained
maximal-severity activity is bounded by the highest severity weight in the
system (0.90 for WORMS; observed maxima of 0.835--0.879 in
Table~\ref{tab:escalation_mode}, consistent with this bound and with the
specific attack-type mixture each intent produces) rather than by 1.0, and
typically sits well below \texttt{RATE\_THRESHOLD} except during genuinely
severe, sustained activity.

\textbf{Interpretation.} This is the empirical confirmation, anticipated in
Section~\ref{sec:escalation_tracking}, that \texttt{RATE\_THRESHOLD}'s
calibration is signal-dependent: a threshold tuned so that R3 "fires only
under sustained, dense attack traffic" (Section~\ref{sec:results_params})
achieves something close to that design intent under EMA tracking, but
under window tracking and continuous-attack conditions, R3 is in practice
the \emph{default} state for most of an episode rather than a rare
escalation trigger. Neither behaviour is incorrect per the respective
signal's definition, but they are materially different operating
characteristics, and Section~\ref{sec:results_adaptive} shows the
consequence of this difference for the bounded threshold controller of
Section~\ref{sec:threshold_controller}. \texttt{escalation\_mode} remains
\texttt{"window"} by default throughout this thesis for continuity with
Section~\ref{sec:results_sedm}--Section~\ref{sec:results_params}'s reported
numbers; adopting EMA tracking operationally would warrant an independent
re-calibration of \texttt{RATE\_THRESHOLD} against the EMA signal's own
statistics, rather than reusing the window-calibrated default, a point now
supported by direct measurement rather than the analytical argument given
in Section~\ref{sec:escalation_tracking} alone.

%% ── 4.11 Cross-Session Reputation and the R4 Override ────────────────
\section{Cross-Session Reputation and the R4 Override}
\label{sec:results_reputation}

This section evaluates the \texttt{ReputationTracker} and R4 override
(Section~\ref{sec:reputation}) using a dedicated simulated-visit protocol,
since — as noted in Section~1.5 — neither the synthetic Gymnasium
environment nor the primary evaluation harness has a cross-episode
source-identity concept to exercise this mechanism directly.

\paragraph{Protocol.}
30 simulated source addresses were each driven through a sequence of up to
8 "offending visits." Each visit consists of one EXPLOITS-severity event
($s_a = 0.70$, Table~\ref{tab:attack_severity_meth}) followed by forced
session expiry (simulating the source returning after its session TTL has
lapsed — reputation, stored independently of the session, persists across
this expiry by design, Section~\ref{sec:reputation}). After each offending
visit, the same source's \emph{next} visit is simulated opening with a
single RECONNAISSANCE-stage, NORMAL-labelled event — the most
innocuous-looking event the protocol can generate — and the action
\texttt{MatrixPolicy} selects for that opening event is compared under two
conditions: \texttt{reputation} supplied from the tracker's actual current
value for that source ("R4 active"), and \texttt{reputation} forced to 0.0
(the pre-R4 baseline, in which R1 unconditionally allows this event).

\begin{table}[ht]
  \centering
  \caption{Mean reputation score and fraction of the 30 simulated sources
           for which R4 overrides R1's opening-event decision, as a
           function of prior offending-visit count.}
  \label{tab:reputation}
  \begin{tabular}{ccc}
    \toprule
    Offending visits & Mean reputation & R4 overrides R1 \\
    \midrule
    1 & 0.175 & 0.0\%             \\
    2 & 0.350 & 0.0\%             \\
    3 & 0.525 & 0.0\%             \\
    4 & 0.700 & \textbf{100.0\%}  \\
    5 & 0.875 & 100.0\%           \\
    6 & 1.000 & 100.0\%           \\
    7 & 1.000 & 100.0\%           \\
    8 & 1.000 & 100.0\%           \\
    \bottomrule
  \end{tabular}
\end{table}

\begin{figure}[ht]
  \centering
  \includegraphics[width=0.82\textwidth]{reputation_threshold_crossing}
  \caption{Cross-session reputation growth (left axis) and the fraction of
           simulated sources for which R4 overrides R1 (right axis), as a
           function of prior offending-visit count.
           The threshold crossing at 4 visits is exact and identical across
           all 30 simulated sources.}
  \label{fig:reputation_threshold_crossing}
\end{figure}

\textbf{A clean, deterministic threshold crossing.} Reputation grows by
exactly $\Delta \cdot s_a = 0.25 \times 0.70 = 0.175$ per offending visit
(decay is negligible at the wall-clock timescale of this simulation, so
growth is effectively linear), crossing \texttt{REPUTATION\_THRESHOLD}
$= 0.60$ at precisely the fourth offending visit (score 0.700). At that
visit and every subsequent one, R4 overrides R1 for all 30 simulated
sources, with no variance across sources — the mechanism's behaviour is, as
intended by its design (Section~\ref{sec:reputation}), fully determined by
the stated arithmetic rather than by any stochastic component.

\textbf{Interpretation.} This result is the direct, quantitative
counterpart to the qualitative trade-off stated in
Section~\ref{sec:reputation}: after a specific, auditable number of prior
offenses (four, for an EXPLOITS-severity attacker under the default
parameters — fewer for a higher-severity attacker such as WORMS, more for a
lower-severity one such as RECONNAISSANCE, since the increment scales with
severity), a returning source's superficially benign opening event is no
longer treated as a fresh first contact. The threshold at which this
occurs, and its sensitivity to the offense severity and the
\texttt{offense\_increment}/\texttt{REPUTATION\_THRESHOLD} constants, is now
known precisely rather than only qualitatively (Section~5.4 revisits
whether these specific default values are appropriate for a real
deployment). This directly targets the "camouflaging attacker" concern
discussed and then withdrawn as an evaluation artefact in
Section~\ref{sec:results_extended}: while that specific empirical claim did
not survive the Bug 8/9 correction, the underlying concern — a source
alternating attack and benign-looking traffic across multiple sessions to
build trust — is exactly the scenario R4 is designed to address, and
Table~\ref{tab:reputation} demonstrates that the mechanism does so, on this
simulated protocol, deterministically and after a small, known number of
offenses.

%% ── 4.12 Bounded Threshold Adaptation ─────────────────────────────────
\section{Bounded Threshold Adaptation}
\label{sec:results_adaptive}

This section evaluates \texttt{AdaptiveThresholds}
(Section~\ref{sec:threshold_controller}) by driving a long decision sequence
through \texttt{MatrixPolicy} twice — once with \texttt{RATE\_THRESHOLD}
fixed at its default, once with an \texttt{AdaptiveThresholds} instance
attached (\texttt{target\_rate}$= 0.10$, \texttt{observation\_window}$= 200$)
— using the identical underlying attacker trajectory (AGGRESSIVE intent,
window-mode escalation tracking, same seed) so the two runs are directly
comparable. 6\,000 decisions (30 observation windows) were recorded under
two traffic regimes: \textbf{continuous attack} (\texttt{benign\_ratio}$=0$,
the regime characterised in Section~\ref{sec:results_escalation_mode} as
having a window-mode R3 rate far above any reasonable target) and
\textbf{mixed traffic} (\texttt{benign\_ratio}$=0.30$, matching
Section~\ref{sec:results_extended}'s extended-evaluation protocol).

\begin{table}[ht]
  \centering
  \caption{Mean R3 trigger rate over the final five observation windows,
           static vs.\ adaptive \texttt{RATE\_THRESHOLD}, under two traffic
           regimes.}
  \label{tab:adaptive_thresholds}
  \begin{tabular}{lccc}
    \toprule
    Regime & Static rate & Adaptive rate & Final threshold \\
    \midrule
    Continuous attack (\texttt{benign\_ratio}$=0$)    & 45.70\% & 44.40\% & 0.90 (bound) \\
    Mixed traffic (\texttt{benign\_ratio}$=0.30$)     & 7.60\%  & 4.40\%  & 0.90 (bound) \\
    \bottomrule
  \end{tabular}
\end{table}

\begin{figure}[ht]
  \centering
  \includegraphics[width=0.88\textwidth]{adaptive_thresholds_convergence}
  \caption{R3 trigger rate per observation window, static vs.\ adaptive
           \texttt{RATE\_THRESHOLD}, under continuous-attack and
           mixed-traffic regimes. The controller saturates at its upper
           bound in both regimes but its effect on the realised trigger
           rate differs sharply between them.}
  \label{fig:adaptive_thresholds_convergence}
\end{figure}

\textbf{The controller correctly identifies excessive firing and responds —
but its bounded design has real, regime-dependent limits.} In both regimes
the controller moves \texttt{RATE\_THRESHOLD} immediately to its configured
upper bound (0.90, i.e.\ \texttt{initial\_threshold} + \texttt{bound}) and
holds it there, correctly diagnosing that R3 is firing far more often than
the 10\% target in both cases. The \emph{consequence} of that response,
however, differs sharply between the two regimes. Under continuous attack,
where Section~\ref{sec:results_escalation_mode} already established that
the window-mode escalation rate saturates near 1.0 for most of an episode,
raising the threshold from 0.80 to 0.90 barely moves the realised trigger
rate (45.70\% $\to$ 44.40\%): the escalation signal is so far above even
the raised bar that the bound is simply insufficient to bring the rate
anywhere near the 10\% target. Under mixed traffic, where
Section~\ref{sec:results_extended}'s 30\% benign injection gives the
escalation signal materially more room below saturation, the same bounded
response is considerably more effective, pulling an already-closer-to-target
static rate of 7.60\% down to 4.40\%.

\textbf{Interpretation.} This is reported as a genuine, and mildly
humbling, empirical finding rather than smoothed over: a bounded
controller, by design (Section~\ref{sec:threshold_controller},
Section~\ref{sec:bg_repctrl}), cannot fully compensate for an escalation
signal that is saturated for structural reasons unrelated to the
controller itself. The result does not undermine the controller's stated
purpose (Section~\ref{sec:threshold_controller}) — it was never framed as a
correctness mechanism, only as an alert-fatigue safety valve — but it does
show that the \emph{bound} chosen ($\pm 0.10$ from the initial threshold)
is calibrated for realistic, non-saturated operating conditions such as the
mixed-traffic regime, and is of limited practical use under the
fully-saturated continuous-attack regime characterised in
Section~\ref{sec:results_escalation_mode}. This is exactly the kind of
finding the controller's honest scoping in
Section~\ref{sec:threshold_controller} was designed to surface rather than
obscure: had the controller instead been framed as a general correctness
mechanism, this result would read as a failure; framed correctly as a
bounded, narrowly-scoped safety valve, it reads as the mechanism behaving
exactly as specified, with a clearly identified operating regime in which
its bound is the binding constraint. Section~5.4 discusses whether the
bound itself, or the choice to calibrate \texttt{AdaptiveThresholds}
against the window signal rather than the EMA signal characterised in
Section~\ref{sec:results_escalation_mode} as less prone to saturation in
the first place, would be the more appropriate point of adjustment.
